% T1 — Ejemplo: procesamiento de una encuesta
% Compilar SIEMPRE en build/:
%   pdflatex -output-directory=build main.tex
\documentclass[11pt,a4paper]{article}

\usepackage[utf8]{inputenc}
\usepackage[T1]{fontenc}
\usepackage[spanish,es-nolists]{babel}
\usepackage[margin=2.5cm]{geometry}
\usepackage{graphicx}
\usepackage{amsmath,amssymb}
\usepackage{booktabs}
\usepackage[hidelinks]{hyperref}

% Los gráficos generados por los scripts viven en ../resultados
\graphicspath{{../resultados/}}

\title{\textbf{T1 — Ejemplo: procesamiento de una encuesta}}
\author{[Nombre del estudiante] \\ Metodologías de Investigación Aplicada (13854.0 M-40)}
\date{\today}

\begin{document}
\maketitle

\begin{abstract}
Ejemplo del flujo completo de una tarea: datos crudos en \texttt{datos/},
procesamiento con un script Python y tablas de resultados para el informe.
\end{abstract}

\section{Introducción}
Esta carpeta es una plantilla de referencia. Al crear la tarea real, se copia
esta estructura y se reemplazan los datos y scripts de ejemplo.

\section{Metodología}
El script \texttt{scripts/procesar\_encuesta.py} lee
\texttt{datos/encuesta.csv} (solo lectura) con el helper
\texttt{utilidades/leer\_excel.py}, calcula estadísticas descriptivas usando
solo la biblioteca estándar y escribe tablas limpias en \texttt{resultados/}.

\section{Resultados}

La Tabla~\ref{tab:genero} muestra la frecuencia por género generada por el
script (ver \texttt{resultados/frecuencia\_por\_genero.csv}).

\begin{table}[htbp]
    \centering
    \begin{tabular}{lcc}
        \toprule
        Género & n & \% \\
        \midrule
        Femenino & 6 & 50.0 \\
        Masculino & 6 & 50.0 \\
        \bottomrule
    \end{tabular}
    \caption{Frecuencia por género (datos de ejemplo).}
    \label{tab:genero}
\end{table}

\section{Conclusiones}
El flujo funciona de punta a cabo: datos $\rightarrow$ script $\rightarrow$
resultados $\rightarrow$ informe.

\end{document}
